\chapter{Schlussbetrachtung}

Ziel der Arbeit war die Entwicklung und Evaluation eines auf Basis der PSI-Theorie fundierten KI-basierten Coaching-Ansatzes. Dabei wurde die Frage betrachtet, wie psychologische Zustandsanalysen mit einer interventionsbasierten Dialogsteuerung verbunden werden können. Im Mittelpunkt stand damit die Übertragbarkeit psychologischer Theorieannahmen auf eine technische Systemarchitektur.

Die Ergebnisse des Versuchs haben gezeigt, dass die PSI-basierte Steuerungslogik bezüglich der ersten Modulationsannahme (Willensbahnung) nach \cite{Kuhl2000The} einen signifikanten zusätzlichen Effekt von +9,99 Prozentpunkte gegenüber der Kontrollgruppe erzielte.
In Betrachtung der zweiten Modulationsannahme zur Selbstentwicklung durch Affektregulation konnte jedoch keine Überlegenheit nachgewiesen werden. Die PSI-Steuerungslogik weiste -0,54 Prozentpunkte im direkten Vergleich zum konventionellen Modell Gemini-3.8-flash auf.
Die allgemeine traumasensible Kontrollbedingung führte in beiden Szenarien ebenso zu positiven Verschiebungen. Die daraus resultierende Erkentnis ist, dass die PSI-Steuerung nicht pauschal überlegen ist. 

%Forschungsbeitag
Die Arbeit liefert einen wissenschaftlichen Beitrag, indem dargestellt wird, wie Annahmen der PSI-Theorie in eine regelbasierte Entscheidungslogik übersetzt werden können.
Durch die Verbindung von Zustandsanalyse, Zuordnung psychischer Funktionssysteme und interventionsbasierte Dialogsteuerung wird sichtbar, wie sich psychologische Prozesse unterschiedlich gut sprachlich operationalisieren lassen. Der Beitrag liegt daher nicht im Nachweis einer psychologischen oder therapeutischen Wirksamkeit, sondern in der konzeptionellen und technischen Umsetzung.
Es wurde ein integriertes Framework entwickelt, das Zustandsanalyse, PSI-basierte Entscheidungslogik, Interventionssteuerung und Sicherheitsregeln miteinander verbindet. Dies geht über eine reine Prompt-Gestaltung und liefert Anknüfungspunkte für zukünftige Forschungen.
Der empirische Betrag besteht in der Untersuchung von 400 Simulationsläufen, die zeigen, wie sich die beiden PSI-Modulationsannahmen durch eine zusätzliche regelbasierte Steuerungslogik sprachlich operationalisieren lassen. Die unterschiedlichen Ergebnisse für M1 und M2 illustrieren, dass der Forschungsbeitrag nicht darin besteht, zu beweisen, dass der Chatbot Menschen bei der Traumaüberwindung hilft. Es wird dargestellt, dass sich psychologische Theorieannahmen unterschiedlich gut in eine technische Dialogsteuerung übersetzen lassen.

Die entwicklete PSI-Steuerungslogik kann insbesondere bei klar operationalisierbaren Interventionszielen einen zusätzlichen Nutzen bieten, wobei komplexere Systemarchitekturen nicht automatisch besser als gut gestaltete Standard-Prompts sind.
Die Systemkomplexität sollte in einem angemessenen Verhältnis zum erwartbaren zusätzlichen Nutzen betrachtet werden. Hinsichtlich des Ziels, eine traumasensible Unterstützung zu gewährleisten, muss eine traumasensible Dialogführung und eine mehrstufige Safety-Architektur zentrale Gestaltungsanforderungen gelten.

Durch das szenariobasierte experimentelle Design zeigten sich einige Limitationen, die in der Arbeit kritisch betrachtet wurden. Es wurden lediglich zwei simulierte Szenarien untersucht, die nur einen begrenzten Ausschnitt möglicher Coaching-Situationen abbilden.
Der Fokus war, linguistische Indikatoren zu analysieren. Diese erlauben jedoch keine unmittelbaren Rückschlüsse auf tatsächliche psychische Zustände, was die Interpretation der Ergebnisse erschwert.
Es wurden keine realen Nutzerinnen und Nutzer und keine klinischen Vorher-Nachher-Veränderungen, sondern nur eine theoriegeleitete Simulation untersucht. Die resultierenden Ergebnisse beziehen sich auf bestimmte Modell- und Prompt-Konfiguration. Obwohl die Resultate Hypothesen aufdecken, konnte keine Konkrete Erklärungen für die Unterschiede zwischen den Bedingungen nicht eindeutig identifiziert werden.

Zusammenfassend lässt sich festhalten, dass das theoretisch-konzeptionelle Ziel der Arbeit erreicht wurde, indem zentrale Annahmen der PSI-Theorie in ein strukturiertes Modell für die Gestaltung und Steuerung eines KI-basierten Coaching-Chatbots überführt wurden. Die entwickelten Zustands-, Ziel- und Interventionslogiken konnten zudem prototypisch technisch umgesetzt und anhand kontrollierter Dialogsimulationen untersucht werden. Damit konnte die grundsätzliche technische Umsetzbarkeit einer theoriegeleiteten Dialogsteuerung demonstriert werden. Die empirische Überprüfung der angestrebten traumasensiblen Unterstützung ist hingegen differenziert zu bewerten. Während für die Modulationsannahme der Willensbahnung ein signifikanter zusätzlicher sprachlicher Steuerungseffekt gegenüber der Kontrollbedingung festgestellt wurde, konnte für die Selbstentwicklung durch Affektregulation keine entsprechende Überlegenheit nachgewiesen werden. Da die Untersuchung ausschließlich auf simulierten Dialogen und linguistischen Indikatoren basiert, lassen sich daraus weder tatsächliche psychologische Veränderungen noch eine therapeutische Wirksamkeit ableiten. Die Arbeit leistet somit einen konzeptionellen und technischen Beitrag zur Übertragung psychologischer Wirkannahmen in KI-basierte Coaching-Systeme, während die tatsächliche Wirksamkeit und Eignung für eine traumasensible Unterstützung in realen Coaching-Interaktionen Gegenstand weiterführender empirischer Untersuchungen bleiben.


\section{Beantwortung der Forschungsfrage}

Für die Beantwortung der Forschungsfrage ergibt sich daraus, dass eine gezielte Verknüpfung von Zustandsanalyse, affektiver Bewertung und interventionsbasierter Dialogsteuerung ein relevantes Designmerkmal für KI-basierten Coaching-Systeme darstellt. Die Ergebnisse sprechen dafür, dass die explizite Berücksichtigung psychologischer Zustände eine differenziertere und zielgerichtetere Interaktion ermöglicht als eine ausschließlich reaktive Gesprächsführung.
Gleichzeitig muss berücksichtigt werden, dass die Untersuchung auf kontrollierten Dialogsimulationen und einer linguistischen Operationalisierung psychischer Systeme basiert. Die gemessenen Systemverschiebungen stellen daher keine unmittelbaren Veränderungen tatsächlicher psychischer Zustände von Personen dar. Insbesondere kann aus der Aktivierung sprachlicher Indikatoren für das Extensionsgedächtnis oder die intuitive Verhaltenssteuerung nicht unmittelbar auf eine tatsächliche psychologische Veränderung geschlossen werden.
Zusammenfassend lässt sich festhalten, dass die Evaluation die Funktionsfähigkeit des entwickelten Frameworks teilweise bestätigt. Für die erste Modulationsannahme konnte ein starker Steuerungseffekt nachgewiesen werden. Für die zweite Modulationsannahme konnte die angenommene Überlegenheit hingegen nicht bestätigt werden. Das entwickelte Modell stellt damit einen konzeptionellen Ansatz dar, dessen Potenzial insbesondere hinsichtlich der Willensbahnung durch die vorliegenden Simulationsergebnisse gestützt wird. Eine weiterführende Untersuchung ist erforderlich, um die Übertragbarkeit der Ergebnisse auf reale Coaching-Interaktionen und die tatsächliche Unterstützung selbstregulatorischer Prozesse zu überprüfen.
Für die Beantwortung der Forschungsfrage lässt sich festhalten, dass die Übertragung der PSI-Theorie auf KI-basierte Coaching-Systeme insbesondere durch die Verknüpfung von fünf zentralen Design- und Interaktionsmerkmalen erfolgen kann. Erstens bildet eine kontextabhängige Zustandsanalyse die Grundlage, um psychologische und affektive Indikatoren über den Gesprächsverlauf hinweg zu erfassen und einzuordnen. Zweitens ermöglicht eine PSI-basierte Entscheidungslogik, die ermittelten sprachlichen Muster den theoretisch relevanten psychischen Funktionssystemen zuzuordnen. Drittens kann eine adaptive Interventionssteuerung darauf aufbauend geeignete sprachliche Impulse auswählen und an der jeweiligen Ausgangslage sowie dem angestrebten Funktionsziel ausrichten. Viertens sollte die Interaktion durch traumasensible Gestaltungsprinzipien geprägt sein, die Ressourcenorientierung, emotionale Sicherheit und die Vermeidung einer Überforderung berücksichtigen. Fünftens ist eine mehrstufige Safety-Architektur notwendig, die Risiken innerhalb der Interaktion erkennt, die Sicherheit der generierten Ausgaben berücksichtigt, transparente Systemgrenzen gewährleistet und bei Bedarf eine menschliche Eskalation ermöglicht.
Die vorliegende Arbeit zeigt, dass sich insbesondere die ersten drei Merkmale in Form einer promptbasierten, PSI-theoretisch fundierten Entscheidungs- und Dialogsteuerung prototypisch umsetzen lassen. Die Evaluation bestätigt dabei einen deutlichen zusätzlichen sprachlichen Steuerungseffekt für die erste Modulationsannahme der Willensbahnung. Für die zweite Modulationsannahme der Selbstentwicklung durch Affektregulation konnte hingegen keine signifikante Überlegenheit gegenüber der bereits traumasensibel gestalteten Kontrollbedingung festgestellt werden. Die Ergebnisse stützen somit die technische Umsetzbarkeit einer theoriegeleiteten Dialogsteuerung, weisen jedoch zugleich auf Grenzen ihrer Übertragbarkeit auf unterschiedliche psychische Prozesse hin. Da die Untersuchung ausschließlich auf simulierten Dialogen und linguistischen Indikatoren beruht, kann daraus keine tatsächliche Verbesserung menschlicher Selbstregulation oder eine therapeutische Wirksamkeit abgeleitet werden. Die fünf genannten Merkmale sind daher als konzeptionelle Gestaltungsanforderungen zu verstehen, deren Nutzen und Wirksamkeit in weiterführenden Untersuchungen mit realen Nutzenden und geeigneten psychologischen Messverfahren überprüft werden müssen.


\section{Ausblick}
Aus den gewonnenen Erkentnissen können Ansätze für zukünftige Forschungen gewonnen werden. 
Es wäre interessant, einen ähnlichen Versuchsaufbau mit realen Testpersonen und vielfältigeren Gesprächssituationen zu testen, um entstandene Hypothesen zu adressieren.
Die Analyse der linguistischen Indikatoren könnte durch geeignete psychologische Messverfahren ergänzt werden sowie eine präzisere Operationalisierung der Interventionsregeln überlegt werden. Vor allem für die zweite Modulationsannahme besteht hier noch Entwicklungspotenzial. Möglicherweise wäre eine Zusammenarbeit mit Experten oder Forschenden im Bereich der Psychologie ein Mehrwert. In diesem Rahmen könnte auch unterscuht werden, ob körperbezogene oder multimodale Interaktionsformen die Unterstützung komplexerer Selbstentwicklungsprozesse verbessern können.
Ein spannender inhaltlicher Aspekt ist außerdem die Weiterentwicklung der Safety-Architektur einschließlich der Prüfung generierter Antworten.
Des Weiteren könnten kontrollierte Vergleichsversuche durchgeführt werden, um Replikation mit unterschiedlichen Sprachmodellen und Modellkonfigurationen zu testen.

